\subsection{A lifting theorem}

PROPOSITION 1. --- Let \(k_{0}\) be a field, A a \(k_{0}\) -algebra which is a separated and complete local ring, and K a sub- \(k_{0}\) -extension of \(\kappa_{\mathrm{A}}\) which possesses a separating transcendence basis \((\xi_{\lambda})_{\lambda \in \Lambda}\) over \(k_{0}\) (A, V, p. 130, def. 1). For every \(\lambda \in \Lambda\) , let \(x_{\lambda}\) be a representative of \(\xi_{\lambda}\) in A. There exists a unique subfield L of A, containing \(k_{0}\) and the elements \(x_{\lambda}\) , such that the canonical homomorphism \(\pi\) from A onto \(\kappa_{\mathrm{A}}\) induces an isomorphism of L onto K.

Let \(\varphi\) be the \(k_{0}\) -homomorphism from the polynomial ring \(k_{0}[(X_{\lambda})_{\lambda \in \Lambda}]\) to A which maps \(X_{\lambda}\) to \(x_{\lambda}\) for every \(\lambda \in \Lambda\) . Let \(u\) be a nonzero element of \(k_{0}[(X_{\lambda})_{\lambda \in \Lambda}]\) ; one has \(\pi (\varphi (u)) \neq 0\) , because the family \((\xi_{\lambda})_{\lambda \in \Lambda}\) is algebraically free over \(k_{0}\) in \(\kappa_{\mathbf{A}}\) ; consequently, \(\varphi (u)\) is invertible in the local ring A. It follows that \(\varphi\) extends to a homomorphism \(\psi\) from the field \(k_{1} = k_{0}((X_{\lambda})_{\lambda \in \Lambda})\) to A. Then A is a \(k_{1}\) -algebra, \(\kappa_{\mathbf{A}}\) is an extension of \(k_{1}\) and K a subextension of \(\kappa_{\mathbf{A}}\) that is algebraic and separable over \(k_{1}\) . It remains to prove that there exists a unique subfield L of A containing \(\psi(k_{1})\) and such that \(\pi(L) = K\) .

\begin{enumerate}
\def\labelenumi{\alph{enumi})}
\item
  Existence of L: Let S be the set of subfields L of A containing \(\psi(k_1)\) and such that \(\pi(L) \subset K\) ; it is inductive for the inclusion relation. Let L be a maximal element of S; we consider K as an algebraic and separable extension of L (A, V, p. 40, Prop. 9). Let \(\xi \in K\) and let \(P \in L[X]\) be its minimal polynomial over L. Since \(\xi\) is a simple root of P, Hensel's lemma (III, § 4, no. 5, cor. 1 of Th. 2) guarantees the existence of an element \(x\) of A such that \(\pi(x) = \xi\) and \(P(x) = 0\) . The subring \(L[X]\) of A belongs to S; by the maximality of L, we therefore have \(x \in L\) , whence \(\xi \in \pi(L)\) . Finally, we have \(\pi(L) = K\) .
\item
  Uniqueness of L: Let L and L' be two subfields of A containing \(\psi(k_1)\) and such that \(\pi(L) = \pi(L') = K\) . Let \(\xi \in K\) , and let \(x \in L\) and \(x' \in L'\) be the elements such that \(\pi(x) = \pi(x') = \xi\) . If \(P \in k_1[X]\) is the minimal polynomial of \(\xi\) over \(k_1\) , then \(\xi\) is a simple root of P, and we have \(P(x) = P(x') = 0\) . By Hensel's lemma (loc. cit.) we have \(x = x'\) . One therefore has \(L = L'\) .
\end{enumerate}

Remark. --- * The preceding proof applies more generally to the case where one assumes only that A is a Henselian local ring*. The uniqueness proof uses the hypothesis that the local ring A is separated, but not that it is complete.

